% LaTeX template for manuscript methods section
% Comprehensive methods section for CFD publication

\section{Methods}

\subsection{Computational Framework}

\subsubsection{Geometry Processing and Mesh Generation}
Patient-specific aortic geometries were obtained from {{ study_metadata.imaging_modality | default('contrast-enhanced computed tomography (CT)') }} with slice thickness of {{ study_metadata.slice_thickness | default('0.6') }} mm. Geometry reconstruction and surface mesh generation were performed using established protocols \cite{ref_geometry_processing}. The computational domain was discretized using the snappyHexMesh utility in OpenFOAM 12 \cite{openfoam_user_guide}, employing a multi-resolution hexahedral mesh with boundary layer refinement near vessel walls.

Mesh independence was verified through Richardson extrapolation using three systematically refined meshes with refinement ratio r = $\sqrt{2}$. The finest mesh contained approximately {{ simulation_config.mesh.finest_cell_count | default('2.5 × 10^6') }} cells, ensuring y+ < 1 in the boundary layer to resolve viscous sublayer dynamics accurately.

\subsubsection{Governing Equations and Numerical Methods}
Blood flow was modeled as an incompressible, Newtonian fluid with density $\rho$ = 1060 kg/m³ and dynamic viscosity $\mu$ = 0.004 Pa·s, consistent with physiological values at body temperature \cite{ref_blood_properties}. The governing equations are the continuity and Navier-Stokes equations:

\begin{equation}
\nabla \cdot \mathbf{u} = 0
\end{equation}

\begin{equation}
\frac{\partial \mathbf{u}}{\partial t} + (\mathbf{u} \cdot \nabla)\mathbf{u} = -\frac{1}{\rho}\nabla p + \nu \nabla^2 \mathbf{u}
\end{equation}

where $\mathbf{u}$ is the velocity vector, $p$ is pressure, and $\nu$ is kinematic viscosity.

Temporal discretization employed a second-order implicit scheme with adaptive time stepping to maintain Courant number < 0.5. Spatial discretization used second-order finite volume schemes: linear interpolation for diffusion terms, linearUpwind for convection terms, and PISO pressure-velocity coupling with 2 corrector steps.

\subsection{Boundary Conditions}

\subsubsection{Inlet Conditions}
Time-varying velocity profiles were prescribed at the aortic root based on patient-specific flow measurements from {{ study_metadata.inlet_data_source | default('phase-contrast MRI') }}. The velocity profile was spatially distributed using a Womersley solution for pulsatile pipe flow \cite{ref_womersley}, accounting for the entrance length development.

\subsubsection{Outlet Boundary Conditions}
Physiologically realistic outlet conditions were implemented using three-element Windkessel (3EWK) models \cite{ref_windkessel}. The 3EWK parameters were determined using Murray's law for automatic flow distribution:

\begin{equation}
\frac{Q_i}{Q_{total}} = \frac{r_i^3}{\sum_{j=1}^{N} r_j^3}
\end{equation}

where $Q_i$ is the flow rate through outlet $i$, $r_i$ is the equivalent radius, and $N$ is the number of outlets.

The 3EWK model relates pressure and flow rate through:

\begin{equation}
p(t) = Z_c Q(t) + R \int_0^t Q(\tau) d\tau + \frac{1}{C} \int_0^t [Q(\tau) - Q_{mean}] d\tau
\end{equation}

where $Z_c$ is characteristic impedance, $R$ is peripheral resistance, and $C$ is arterial compliance.

\subsubsection{Wall Conditions}
Vessel walls were treated as rigid with no-slip boundary conditions. This assumption is justified for aortic flow analysis where wall motion effects are secondary to hemodynamic patterns \cite{ref_rigid_wall_assumption}.

\subsection{Solution Procedure and Convergence Criteria}

Simulations were performed using the incompressibleFluid solver in OpenFOAM 12. Each cardiac cycle (T = {{ simulation_config.cardiac_cycle_period | default('0.8') }} s, corresponding to {{ (60 / (simulation_config.cardiac_cycle_period | default(0.8))) | round }} bpm) was discretized with 2000 time steps, ensuring temporal resolution of peak flow gradients.

Convergence criteria were set to:
\begin{itemize}
    \item Velocity residuals: < $10^{-6}$
    \item Pressure residuals: < $10^{-7}$
    \item Mass flux imbalance: < 0.1\%
\end{itemize}

Simulations were run for 4 cardiac cycles to eliminate transient effects, with data analysis performed on the final cycle.

\subsection{Hemodynamic Analysis}

\subsubsection{Wall Shear Stress Metrics}
Time-averaged wall shear stress (TAWSS) was calculated as:

\begin{equation}
\text{TAWSS} = \frac{1}{T} \int_0^T |\boldsymbol{\tau}_w(t)| dt
\end{equation}

Oscillatory shear index (OSI) was computed as:

\begin{equation}
\text{OSI} = \frac{1}{2}\left(1 - \frac{\left|\int_0^T \boldsymbol{\tau}_w(t) dt\right|}{\int_0^T |\boldsymbol{\tau}_w(t)| dt}\right)
\end{equation}

Relative residence time (RRT) was defined as:

\begin{equation}
\text{RRT} = \frac{1}{(1 - 2 \cdot \text{OSI}) \cdot \text{TAWSS}}
\end{equation}

\subsubsection{Energy Loss Analysis}
Viscous energy loss was quantified using the energy dissipation rate:

\begin{equation}
\Phi = \mu \left(\frac{\partial u_i}{\partial x_j} + \frac{\partial u_j}{\partial x_i}\right)^2
\end{equation}

Pressure drop was calculated between inlet and outlets, accounting for hydrostatic effects.

\subsection{Validation and Verification}

\subsubsection{Verification Studies}
Grid convergence index (GCI) was calculated using Richardson extrapolation \cite{ref_gci}:

\begin{equation}
\text{GCI} = \frac{1.25 \epsilon}{r^p - 1}
\end{equation}

where $\epsilon$ is the relative error between successive grids, $r$ is the refinement ratio, and $p$ is the observed order of convergence.

Time step independence was verified by halving the time step and ensuring solution changes < 1\%.

\subsubsection{Validation Against Experimental Data}
Computational results were validated against experimental measurements from {{ study_metadata.experimental_setup | default('stereoscopic particle image velocimetry (sPIV) in patient-specific flow phantoms') }}. Validation metrics included:

\begin{itemize}
    \item Pearson correlation coefficient (R²)
    \item Root mean square error (RMSE)
    \item Bland-Altman limits of agreement
    \item Concordance correlation coefficient
\end{itemize}

\subsubsection{Literature Comparison}
Results were compared with published values from similar patient populations, using statistical tests to assess agreement within reported confidence intervals.

\subsection{Statistical Analysis}
All statistical analyses were performed using Python SciPy \cite{ref_scipy}. Continuous variables are reported as mean ± standard deviation unless otherwise specified. Significance was set at p < 0.05 for all tests.

Uncertainty quantification was performed using Monte Carlo simulation with 1000 samples, varying key parameters within physiologically relevant ranges:
\begin{itemize}
    \item Blood viscosity: 0.003-0.005 Pa·s
    \item Density: 1040-1080 kg/m³
    \item Flow rate amplitude: ±10\% of measured values
\end{itemize}

\subsection{Computational Resources}
All simulations were performed on {{ study_metadata.computational_resources | default('a high-performance computing cluster with Intel Xeon processors') }}. Typical wall clock time for a single cardiac cycle simulation was {{ study_metadata.simulation_time | default('6-8') }} hours using {{ study_metadata.cpu_cores | default('32') }} CPU cores with MPI parallelization.